\subsection{The Bazzi--Mahdian--Spielman Question}
\label{sec:bms-question}

\paragraph{When Does One Interleaver Suffice?}
Bazzi, Mahdian, and Spielman (BMS) show that a depth-two serial concatenation
cannot be asymptotically good when the outer has constant memory and the
inner has sublinear memory~\cite[Sec.~III-A]{bms09}. Their conclusion asks
whether allowing the outer memory to grow logarithmically with block length
permits an asymptotically good code~\cite[Sec.~IV]{bms09}.
Here, asymptotically good means that both rate and relative minimum
distance stay bounded away from zero. The question concerns the
composition of an outer encoder, one interleaver, and an inner encoder.

\paragraph{The SPIN Progression.}
Accumulator SPIN supplies the basic positive construction in the
block-encoder setting. Its outer blocks have length $\Theta(\log N)$,
and its inner accumulator carries one bit of state. At rate $1/2$,
Theorem~\ref{thm:independent-accumulator-spin} gives relative distance
greater than $0.0037$ with probability $1-o(1)$. Thus a growing inner
state is unnecessary for linear distance once the outer blocks grow
logarithmically. Random SPIN increases the inner memory to
$\Theta(\log N)$ and proves relative distance greater than $0.11002$.
Structured SPIN then returns to constant recursive state: with better
constituents and structured interleaving, it proves relative distance
greater than $0.11$ and linear-time encoding. Both of the latter distance
guarantees are close to the rate-$1/2$ GV value of approximately $0.110028$.

\paragraph{Logarithmic Outer Memory Suffices.}
The block description admits a streaming realization with logarithmic
memory and a short termination. For a rate-$1/2$ outer with block length
$b$, buffer one block while emitting the preceding block, at two output
bits per input bit. Initially the output buffer contains $b$ zeros.
After the message, $b/2$ fixed zero input bits flush the last block.
This produces $0^b\Vert O(x)$ using $O(b+\log N)$ memory, including
buffers and a position counter. Extend the interleaver to send these
startup zeros to the end and apply the original permutation to $O(x)$.
The inner therefore receives $\Pi(O(x))\Vert0^b$. Since each inner is
causal, the first $N$ output bits are exactly the original SPIN codeword.
The added outputs cannot reduce absolute minimum distance. For
$b=\Theta(\log N)$, the rate becomes $1/2-o(1)$ and the relative-distance
guarantee changes by only $o(1)$.

This gives a positive answer in the terminated-encoder interpretation of
the BMS question: logarithmic outer memory and a one-bit accumulator
suffice for an asymptotically good code. BMS mention termination but
omit it from their displayed encoder definition~\cite[Sec.~I-A]{bms09};
the fixed termination inputs are explicit in the realization above.
The structured interleaver is still one permutation between the two
encoders; permutations internal to the outer act only within its blocks.
